\section{Chạy thử trên corpus của sách}
\label{sec:ch07-corpus}

\index{Transformer!đặt cạnh mô hình chương 6}%
Sáu mục trước dựng kiến trúc và đo từng mảnh. Phép kiểm thẳng thắn nhất còn lại
là đặt cả cái vào đúng chỗ mô hình chương~\ref{ch:mot-vector-cho-moi-tu} đã
đứng, giữ nguyên mọi thứ khác, và xem con số đi đâu.

Mọi thứ giữ nguyên nghĩa là mọi thứ: cùng corpus, cùng cách chia tập, cùng 6000
cặp huấn luyện và 300 cặp test, cùng batch 128, cùng Adam ở learning rate 0.005,
cùng clip 5.0, cùng phép \tn{đảo câu nguồn}{source reversal}, cùng giải mã tham
lam, cùng cách chấm \tn{khớp đúng cả câu}{exact match}. Đổi hai thứ, và bảng có
một cột cho mỗi thứ: kiến trúc, và số
epoch.

\begin{measured}{mô hình chương trước và Transformer, cùng recipe}
\begin{minted}{text}
model        d     epochs  params   loss     exact    short    long
attention    32    14      40805    0.0088   1.0000   1.0000   1.0000
transformer  24    14      35845    0.1159   0.4700   0.8581   0.0921
transformer  24    28      35845    0.0148   0.8967   1.0000   0.7961
transformer  24    42      35845    0.0008   0.9967   1.0000   0.9934
transformer  24    56      35845    0.0001   0.9967   1.0000   0.9934

short = source of at most 11 tokens (148 of 300), long = the rest
\end{minted}

Transformer 2 tầng mỗi phía, 4 head, \(\dff = 4\dmodel\). Một seed mỗi bảng.

Đo bằng \code{experiments/ch07\_corpus.py} tại tag \code{ch07}.
\end{measured}

\index{đo lường!hàng đối chứng tái lập chương trước}%
Đọc phần không phải kết quả trước. Hàng đầu là mô hình
chương~\ref{ch:mot-vector-cho-moi-tu} huấn luyện lại ở đây chứ không phải chép
sang, và
nó trùng bảng mục~\ref{sec:ch06-be-rong} ở cả năm số: 40805 tham số, mất mát
0.0088, rồi 1.0000 ba lần. Nếu hàng ấy lệch, mọi chênh lệch phía dưới còn có thể
là chênh lệch của môi trường chạy chứ không phải của kiến trúc.

Bốn hàng Transformer cũng không phải bốn lần huấn luyện. Chúng là \emph{một}
lần, chấm ở bốn thời điểm trên đường đi của chính nó. Cùng trọng số, cùng thứ tự
batch, cùng trạng thái optimizer, nên hàng 14 epoch đúng bằng mô hình mà một
lần chạy 14 epoch cho ra.

\subsection*{Ở recipe của chương trước, kiến trúc mới thua}

\index{Transformer!thua ở 14 epoch}%
Hàng thứ hai là câu trả lời cho câu hỏi tôi bắt đầu mục này với, và nó không
phải câu tôi chờ đợi. Ở đúng 14 epoch mà chương~\ref{ch:mot-vector-cho-moi-tu}
chạy, Transformer được 0.4700 so với 1.0000. Trên câu dài nó được 0.0921.

Đó là một khoảng cách không có cách nào đọc thành thắng lợi, và cũng không phải
chuyện làm tròn. Nên phải hỏi tại sao, và có ba cách giải thích cần tách khỏi
nhau trước khi nói bất cứ điều gì.

Ba phép chẩn đoán dưới đây chạy ở \(\dmodel = 32\) chứ không phải ở
\(\dmodel = 24\) của bảng, vì chúng làm xong trước khi bề rộng của bảng được
chốt. Nên đọc chúng như số của một mô hình rộng hơn ở cùng 14 epoch, mô hình
được 0.6533 chứ không phải 0.4700, và cùng kiểu hỏng. Chúng trả lời câu
\enquote{vì sao thấp}, không trả lời câu \enquote{thấp bao nhiêu}.

\textbf{Không phải overfit.} Chấm mô hình \(\dmodel = 32\) ấy trên 300 cặp lấy
từ chính tập huấn luyện thì được 0.5967, không cao hơn 0.6533 nó được trên tập
test. Mô hình chưa khớp nổi cả dữ liệu nó đã thấy, nên chỗ thiếu không phải chỗ
nó học thuộc quá tay.

\textbf{Không phải mã hỏng.} Đưa cho nó prefix đúng rồi đo độ chính xác từng
token thì được 0.9657 trên toàn tập và 0.9499 trên riêng câu dài. Mô hình gần
đúng ở mức token; cái phạt nó là phép chấm. Khớp đúng cả câu đòi mọi token đều
đúng, và \(0.9499^{19} = 0.3765\) cho câu đích 19 token, so với 0.3421 mà chính
mô hình ấy được ở cột \code{long}. Phần lớn khoảng cách giữa hai cột của nó giải
thích được chỉ bằng phép nhân dồn, không cần giả định gì thêm.

\textbf{Không phải optimizer.} Đây là cách giải thích tôi tin nhất lúc đầu, vì
mục 5.3 của bài có hẳn một lịch warmup và recipe dùng chung thì không. Tôi thử
sáu công thức learning rate ở \(\dmodel = 32\), gồm cả lịch warmup của chính
bài thu nhỏ cho vừa 658 bước huấn luyện. Recipe dùng chung, learning rate 0.005
cố định, cho kết quả \emph{tốt nhất} trong cả sáu. Lịch warmup ở warmup 100 và
400 cho 0.4267 và 0.5000, đều thấp hơn 0.6533. Chạy gấp đôi số epoch ở cùng
recipe ấy cho 0.6033, tức cũng không đưa nó đi đâu. Nên chương này không được
đổ cho optimizer, và chương~\ref{ch:chi-phi-va-song-song} vẫn là chỗ đo warmup
cho tử tế ở quy mô nó có nghĩa.

\subsection*{Nó không chạm trần, nó thiếu bước}

\index{Transformer!ba lần số epoch thì tới nơi}%
Ba hàng còn lại là chỗ bảng này đổi nghĩa. Ở 28 epoch, 0.8967. Ở 42 epoch,
0.9967. Ở 56 epoch, vẫn 0.9967, tức nó đã dừng.

Với 35845 tham số, ít hơn 40805 của mô hình attention.

Phát biểu tôi thấy đứng vững: trên corpus này, hai kiến trúc đi tới cùng một
chỗ, Transformer bằng ít tham số hơn và bằng ba lần số epoch. Cái nó cần không
phải sức chứa, mà là số bước cập nhật.

Một chú ý để không ai đọc bảng thành đường thẳng. Cột \code{loss} giảm đơn điệu
qua bốn mốc, 0.1159 rồi 0.0148 rồi 0.0008 rồi 0.0001, còn cột khớp đúng cả câu
thì không phải lúc nào cũng thế: một lần chạy riêng ở \(\dmodel = 32\) cho
0.6533 ở 14 epoch và 0.6033 ở 28 epoch. Khớp đúng cả câu là tích của xác suất
từng token trên cả câu, nên nó nhảy chứ không đi mượt, và một cột như thế cần
nhiều mốc mới đọc được.

\subsection*{Mô hình sai ở đâu, và câu trả lời không phải \enquote{sai ngẫu nhiên}}

\index{Transformer!lỗi gắn nhầm cụm}%
Con số 0.4700 nói mô hình sai nhiều. Nó không nói sai thế nào, và đọc output ra
thì kiểu lỗi rất đều.

\begin{minted}{text}
src  ... and the big bird wants a black horse
want ... và con chim lớn muốn một con ngựa đen
got  ... và con ngựa lớn muốn một con chim đen

src  a dog carries a small horse and a bird carries the white
     horse in the garden
want một con chó mang một con ngựa nhỏ và một con chim mang
     con ngựa trắng trong vườn
got  một con chó nhỏ mang một con ngựa và một con chim trắng
     mang con ngựa trong vườn
\end{minted}

Ở ví dụ thứ nhất, hai danh từ của hai mệnh đề đổi chỗ cho nhau. Ở ví dụ thứ
hai, hai tính từ gắn vào nhầm \tn{cụm danh từ}{noun phrase}. Từ vựng đúng hết,
số lượng token
đúng hết, cấu trúc câu đúng hết. Cái sai là \emph{cái nào thuộc về cái nào}.

Đó là kiểu lỗi rất riêng, và nó là bản xem trước của một chuyện lớn hơn nhiều.
Mạng hồi quy đọc câu theo thứ tự, nên \enquote{tính từ đứng ngay trước danh từ
này} là chuyện nó có sẵn, không phải học. Với self-attention thì không có gì có
sẵn: mọi vị trí nhìn mọi vị trí như nhau, và quan hệ \enquote{ngay trước} phải
được dựng lại từ \tn{mã hóa vị trí}{positional encoding}, bằng dữ liệu. Sáu
nghìn câu của một văn phạm
hữu hạn gần đủ để dựng, và \enquote{gần đủ} là 0.4700.

Cái phải trả để bỏ recurrence không phải sức chứa và không phải thời gian mỗi
bước. Nó là \tn{thiên kiến quy nạp}{inductive bias}, và
chương~\ref{ch:bias-quy-nap} là chương dành
riêng cho chuyện đó.

\subsection*{Bảng này nói được gì về bài báo}

\index{corpus đồ chơi!giới hạn của bảng chương 7}%
Không nhiều, và giới hạn thì vẫn đúng giới hạn mục~\ref{sec:ch06-be-rong} đã
nêu. Văn phạm hữu hạn, \tn{từ điển}{vocabulary} đóng, không có từ
\tn{ngoài từ điển}{out-of-vocabulary}, không nhập
nhằng, câu dài nhất 23 token, và 6000 cặp. Bài chạy WMT'14 với 4.5 triệu cặp
câu. Một mô hình đạt 0.9967 ở đây nghĩa là nó học thuộc được một văn phạm nhỏ.

Cái bảng này kiểm được là một câu điều kiện: nếu bỏ recurrence mà giữ attention
thì mô hình còn dịch được không, và trả bằng cái gì. Câu trả lời trên corpus
này là còn, và trả bằng số bước. Chiều ấy khớp với chính bài ở chỗ bài nói rất
ít: bài không so số bước với một mô hình hồi quy cùng dữ liệu, bài so thời gian
tường và số FLOP, mà hai thứ đó thắng vì lý do khác hẳn.
